\section*{Capítulo \ref{ch:b1:geological-maps} --- Mapas geológicos}

\begin{solution}{exo:b1:geological-maps:1}
\omterm{def:b1:geological-maps:attitude}{Direção}: \omterm{def:b1:spherical-astronomy:horizontal}{azimute} das linhas horizontais do plano; \omterm{def:b1:geological-maps:attitude}{mergulho}: ângulo entre o plano e a horizontal, medido perpendicularmente à \omterm{def:b1:geological-maps:attitude}{direção}; \omterm{def:b1:geological-maps:attitude}{sentido de mergulho}: \omterm{def:b1:spherical-astronomy:horizontal}{azimute} da linha mais inclinada para baixo. Para $135/40$, a \omterm{def:b1:geological-maps:attitude}{direção} pela regra da mão direita é $135 - 90 = \ang{045}$. Símbolo: uma linha de \omterm{def:b1:geological-maps:attitude}{direção} orientada em N\,45\,L, com um traço do lado sudeste (em direção a \ang{135}) rotulado com 40.
\end{solution}

\begin{solution}{exo:b1:geological-maps:2}
$z = 500 - \tan\ang{30}\,(x\sin\ang{90} + y\cos\ang{90}) = 500 - 0.577\,x$.
Em $(1000, 0)$: $500 - 577 = \qty{-77}{m}$. Em $(0, 1000)$: \qty{500}{m}, pois o plano é horizontal ao longo da sua \omterm{def:b1:geological-maps:attitude}{direção} norte--sul.
\end{solution}

\begin{solution}{exo:b1:geological-maps:3}
Um traço que segue uma curva de nível é um contato horizontal, a \qty{350}{m}. Um traço que corre em linha reta, atravessando colinas e vales, pertence a um contato vertical.
\end{solution}

\begin{solution}{exo:b1:geological-maps:4}
Ao longo do rio, o núcleo do \omterm{def:b1:geological-maps:fold}{anticlinal} mostra o xisto e o granito (sob o arenito); o núcleo do \omterm{def:b1:geological-maps:fold}{sinclinal} é ocupado pelo giz. As unidades mais antigas ficam mais baixas na sequência: onde as camadas se arqueiam para cima e a erosão corta mais fundo, aparecem as mais antigas; onde elas se encurvam para baixo, as mais novas se preservam no meio.
\end{solution}

\begin{solution}{exo:b1:geological-maps:5}
Os dois primeiros pontos estão na mesma altitude: a reta $y = 0$ é uma linha de \omterm{def:b1:geological-maps:attitude}{direção}, logo a \omterm{def:b1:geological-maps:attitude}{direção} é leste--oeste. O contato está \qty{100}{m} mais baixo,
\qty{500}{m} ao norte: ele mergulha para o norte, $\tan\delta = 100/500 = 0.2$,
$\delta \approx \ang{11.3}$; \omterm{def:b1:geological-maps:attitude}{sentido de mergulho} \ang{000}, \omterm{def:b1:geological-maps:attitude}{direção} \ang{270} pela regra da mão direita. Com o teorema: $\overrightarrow{P_1P_2} = (1000, 0, 0)$,
$\overrightarrow{P_1P_3} = (0, 500, -100)$, $\vec n = (0, 100\,000,
500\,000)$, $\cos\delta = 500\,000/509\,902 \approx 0.981$, $\delta \approx
\ang{11.3}$, e $(n_x, n_y) = (0, 100\,000)$ aponta para o norte: o mesmo resultado.
\end{solution}

\begin{solution}{exo:b1:geological-maps:6}
A \omterm{def:b1:geological-maps:attitude}{direção} rio abaixo é o norte, $s = 0.03$ (\ang{1.72}). (a) $a = 0 < s$: V rio acima (sul), ao longo das curvas de nível. (b) $a = -\tan\ang{20} < 0 < s$: V rio acima.
(c) $a = \tan\ang{1} \approx 0.017 < s$: V rio acima. (d) $a = \tan\ang{20}
\approx 0.364 > s$: V rio abaixo (norte).
\end{solution}

\begin{solution}{exo:b1:geological-maps:7}
$\xi_{\text{fundo}} - \xi_{\text{flanco}} = h/(a - s) = 80/(0.364 - 0.03)
\approx \qty{240}{m}$: o traço no flanco fica cerca de \qty{240}{m} rio acima do vértice.
\end{solution}

\begin{solution}{exo:b1:geological-maps:8}
O fechamento do \omterm{def:b1:geological-maps:fold}{anticlinal} aponta na \omterm{def:b1:geological-maps:attitude}{direção} do caimento: o eixo cai para o norte. Com um eixo caindo para o norte, os \omterm{def:b1:geological-maps:geological-map}{afloramentos} de um \omterm{def:b1:geological-maps:fold}{sinclinal} se fecham para o sul.
\end{solution}

\begin{solution}{exo:b1:geological-maps:9}
O lado leste-sudeste mostra unidades mais novas (giz e folhelho, contra unidades mais antigas em frente a elas): esse bloco desceu. Os traços do símbolo da falha estão desse lado, o bloco abatido.
\end{solution}

\begin{solution}{exo:b1:geological-maps:10}
Uma reta comum aos dois planos é perpendicular às duas normais, logo paralela a $\vec n_1 \times \vec n_2$. Para $090/30$: $\vec n_1 = (0.577, 0,
1)$; para $270/30$: $\vec n_2 = (-0.577, 0, 1)$. Então $\vec n_1 \times \vec n_2
= (0\cdot 1 - 1\cdot 0,\ 1\cdot(-0.577) - 0.577\cdot 1,\ 0) = (0, -1.155,
0)$: uma reta horizontal norte--sul. O primeiro plano mergulha para leste, o segundo para oeste: eles mergulham afastando-se da reta, que é a crista de um \omterm{def:b1:geological-maps:fold}{anticlinal} de eixo horizontal.
\end{solution}

\begin{solution}{exo:b1:geological-maps:11}
O traço de um dique vertical é a reta que passa pela origem com \omterm{def:b1:geological-maps:attitude}{direção} \ang{045}, qualquer que seja a altitude: na colina ele está em algum ponto dessa mesma reta. O dique que mergulha é o plano $z = -\tan\ang{60}\,d$, em que $d$ é a distância horizontal à reta de \omterm{def:b1:geological-maps:attitude}{direção} que passa pela origem, contada positivamente para sudeste. Em $z = \qty{300}{m}$: $d = -300/\tan\ang{60}
\approx \qty{-173}{m}$: o traço fica \qty{173}{m} a noroeste da reta de \omterm{def:b1:geological-maps:attitude}{direção} que passa pela origem, do lado de cima do plano.
\end{solution}

\begin{solution}{exo:b1:geological-maps:12}
\[
  \frac{\partial z}{\partial x} \approx \frac{405 - 420}{100} = -0.15, \qquad
  \frac{\partial z}{\partial y} \approx \frac{412 - 420}{100} = -0.08,
\]
logo a declividade é $\arctan\sqrt{0.15^2 + 0.08^2} = \arctan 0.17 \approx
\ang{9.6}$. O terreno desce ao longo de $-\nabla z = (0.15, 0.08)$, cujo \omterm{def:b1:spherical-astronomy:horizontal}{azimute} é
$\arctan(0.15/0.08) \approx \ang{62}$: a célula está voltada para o leste-nordeste.
\end{solution}

\begin{solution}{pb:b1:geological-maps:1}
\textbf{1.} Xisto, granito, arenito, calcário, folhelho, giz, aluvião.

\textbf{2.} De oeste para leste: giz e folhelho, calcário, arenito, depois xisto e granito perto de $x = \qty{3.5}{km}$, arenito, calcário e folhelho de novo, depois giz; aluvião ao longo de todo o fundo do vale.

\textbf{3.} O arenito, com o xisto e o granito sob ele ao longo do rio: as unidades mais antigas no núcleo, um \omterm{def:b1:geological-maps:fold}{anticlinal}.

\textbf{4.} O giz, a unidade mais nova: um \omterm{def:b1:geological-maps:fold}{sinclinal}.

\textbf{5.} A base do arenito é mais alta sobre o eixo do \omterm{def:b1:geological-maps:fold}{anticlinal}, e só ao longo do rio a erosão cortou abaixo dela.

\textbf{6.} Uma discordância: o arenito repousa sobre duas unidades mais antigas diferentes, xisto e granito, e corta o contorno do granito.

\textbf{7.} $435 - 261 = 174$, $432 - 428 = 4$, $512 - 586 = \qty{-74}{m}$.

\textbf{8.} $\overrightarrow{E_1E_2} = (600, 1200, -170)$,
$\overrightarrow{E_1E_3} = (1100, 400, -248)$.

\textbf{9.} O produto vetorial é
\[
  (-229\,600,\ -38\,200,\ -1\,080\,000);
\]
orientado para cima, $\vec n_E = (229\,600,\ 38\,200,\ 1\,080\,000)$.

\textbf{10.} $\|\vec n_E\| \approx 1\,104\,800$; $\cos\delta \approx 0.9776$,
$\delta \approx \ang{12.2}$.

\textbf{11.} \omterm{def:b1:spherical-astronomy:horizontal}{Azimute} de $(229\,600, 38\,200)$: $\arctan(229\,600/38\,200)
\approx \ang{80.6}$; \omterm{def:b1:geological-maps:attitude}{atitude} $081/12$, \omterm{def:b1:geological-maps:attitude}{direção} $\approx \ang{351}$.

\textbf{12.} Sim: o símbolo do flanco leste tem uma linha de \omterm{def:b1:geological-maps:attitude}{direção} quase norte--sul e um traço de \omterm{def:b1:geological-maps:attitude}{mergulho} para leste, rotulado com 12.

\textbf{13.} $524 - 669 = -145$, $421 - 459 = -38$, $319 - 209 =
\qty{110}{m}$.

\textbf{14.} $\overrightarrow{W_1W_2} = (700, 1200, 107)$ e
$\overrightarrow{W_1W_3} = (1100, -600, 255)$. Seu produto vetorial, orientado para cima, é
\[
  (370\,200,\ -60\,800,\ -1\,740\,000) \quad\longrightarrow\quad
  \vec n_W = (-370\,200,\ 60\,800,\ 1\,740\,000).
\]

\textbf{15.} Daí $\cos\delta = 1\,740\,000/\|\vec n_W\| \approx 0.9776$, $\delta
\approx \ang{12.2}$; \omterm{def:b1:geological-maps:attitude}{sentido de mergulho}: \omterm{def:b1:spherical-astronomy:horizontal}{azimute} de $(-370\,200, 60\,800)$,
$\approx \ang{279.3}$. \omterm{def:b1:geological-maps:attitude}{Atitude} $279/12$, \omterm{def:b1:geological-maps:attitude}{direção} $\approx \ang{189}$: o flanco oeste mergulha para oeste-noroeste, o flanco leste para leste-nordeste.

\textbf{16.} $\vec n_E \times \vec n_W = (804 \times 10^6,\ -799\,320 \times
10^6,\ 28\,101 \times 10^6)$ (arredondado).

\textbf{17.} O sentido de descida é $(-804, 799\,320, -28\,101) \times
10^6$: rumo $\arctan(-804/799\,320) \approx \ang{0}$, isto é, norte; caimento
$\arctan\bigl(28\,101/\sqrt{804^2 + 799\,320^2}\bigr) \approx \ang{2.0}$. O
eixo cai \ang{2} para o norte.

\textbf{18.} $8000 \times \tan\ang{2.0} \approx \qty{280}{m}$ mais baixo, perto de
\qty{-50}{m}. O núcleo do \omterm{def:b1:geological-maps:fold}{anticlinal} é mais alto perto do rio e afunda para o norte, de modo que o seu \omterm{def:b1:geological-maps:geological-map}{afloramento} se fecha para o norte, como mostra o mapa.

\textbf{19.} $\overrightarrow{F_1F_2} = (1627, 2500, -276)$ e
$\overrightarrow{F_1F_3} = (3124, 5500, 77)$ dão a normal para cima
\[
  (1\,710\,500,\ -987\,503,\ 1\,138\,500), \qquad
  \delta = \arccos\frac{1\,138\,500}{2\,279\,700} \approx \ang{60},
\]
e um \omterm{def:b1:geological-maps:attitude}{sentido de mergulho} $\approx \ang{120}$. \omterm{def:b1:geological-maps:attitude}{Atitude} $120/60$, \omterm{def:b1:geological-maps:attitude}{direção}
\ang{030}.

\textbf{20.} O \omterm{def:b1:geological-maps:attitude}{mergulho} da falha é para leste-sudeste, de modo que esse lado é o teto, o bloco que fica acima do plano de falha inclinado; ele desceu (unidades mais novas na superfície): uma falha normal.

\textbf{21.} O plano da falha mergulha para leste-sudeste, rio abaixo e muito mais acentuadamente do que desce o fundo do vale: terreno mais baixo encontra o plano mais longe no sentido do \omterm{def:b1:geological-maps:attitude}{mergulho}, de modo que o traço se desvia para leste-sudeste no vale --- a regra do V aplicada ao plano da falha.

\textbf{22.} Intrusão do granito; erosão antes do arenito; deposição do giz; dobramento; falhamento; deposição do aluvião.

\textbf{23.} Depois de 90 milhões de anos atrás (o giz está dobrado) e antes de 2.58 milhões de anos atrás (o aluvião não está falhado); o dobramento veio primeiro, o falhamento depois dele.

\textbf{24.} A falha desce \qty{1405}{m} de F1 (\qty{405}{m}) até
\qty{-1000}{m}; com um \omterm{def:b1:geological-maps:attitude}{mergulho} de \ang{60}, isso corresponde a uma distância horizontal
$1405/\tan\ang{60} \approx \qty{811}{m}$ no sentido do \omterm{def:b1:geological-maps:attitude}{mergulho}. O poço deve ser perfurado
cerca de \textbf{\qty[text-series-to-math=true]{810}{m} a leste-sudeste de F1} (\omterm{def:b1:spherical-astronomy:horizontal}{azimute} \ang{120}).
\end{solution}
